% BHU PYQ -- Professional Exam Template
\documentclass[11pt,a4paper]{article}

\usepackage[a4paper,top=2.5cm,bottom=2.5cm,left=2.8cm,right=2.8cm,headheight=14pt]{geometry}
\usepackage{amsmath,amssymb}
\usepackage[T1]{fontenc}
\usepackage[utf8]{inputenc}
\usepackage{lmodern}
\usepackage{microtype}
\usepackage{fancyhdr}
\usepackage{enumitem}
\usepackage{hyperref}
\usepackage{lastpage}
\usepackage{parskip}

\hypersetup{colorlinks=true,urlcolor=black,linkcolor=black,
  pdftitle={B.Sc. (Hons.) Zoology Sem-II ZOB-201 -- BHU PYQ}}

\pagestyle{fancy}
\fancyhf{}
\renewcommand{\headrulewidth}{0.4pt}
\renewcommand{\footrulewidth}{0.4pt}
\fancyhead[L]{\small\textit{Banaras Hindu University}}
\fancyhead[R]{\small\textit{Zoology Paper: ZOB-201}}
\fancyfoot[C]{\small Page \thepage\ of \pageref{LastPage}}

\newlist{parts}{enumerate}{2}
\setlist[parts,1]{label=(\alph*),leftmargin=*,itemsep=4pt,topsep=3pt}
\setlist[parts,2]{label=(\roman*),leftmargin=*,itemsep=2pt,topsep=2pt}

\newlist{choices}{enumerate}{1}
\setlist[choices,1]{label=(\arabic*),leftmargin=*,itemsep=2pt,topsep=2pt}

\newcommand{\pts}[1]{\hfill\textit{[#1]}}

\begin{document}

\begin{center}
  {\large\bfseries BANARAS HINDU UNIVERSITY}\\[2pt]
  {\normalsize B.Sc. (Hons.) Semester II Examination, 2023-24}\\[6pt]
  \rule{0.8\linewidth}{0.5pt}\\[6pt]
  {\large\bfseries Zoology}\\[4pt]
  {\large\bfseries Paper: ZOB-201 --- Fundamental of Cell Biology and Elementary Biochemistry}\\[4pt]
  {\normalsize Time: Three hours \hfill Full Marks: 70}
\end{center}

\rule{\linewidth}{0.4pt}

\smallskip
\noindent\textit{Note: Answer three questions from Section-A (including Question No. 1 which is compulsory) and three questions from Section-B (including Question No. 1 which is compulsory). Use a separate answer book for each section. The figures in the right-hand margin indicate marks.}

\bigskip
\begin{center}
  \textbf{SECTION -- A} \\
  \textbf{(Fundamental of Cell Biology) (Marks: 35)}
\end{center}

\noindent\textbf{Question 1.} Answer the following parts:
\begin{parts}
  \item Select and write the best options in the following:\pts{$4 \times \frac{1}{2} = 2$}
  \begin{parts}
    \item The tails of the phospholipids of the plasma membrane are composed of \rule{1.5cm}{0.4pt} and \rule{1.5cm}{0.4pt}.
    \begin{choices}
      \item phosphate groups, hydrophobic
      \item fatty acid groups, hydrophilic
      \item phosphate groups, hydrophilic
      \item fatty acid groups, hydrophobic
    \end{choices}
    
    \item Which of the following do not participate in innate immunity?
    \begin{choices}
      \item Natural killer cells
      \item Phagocytic cells
      \item Polymorphonuclear leukocytes
      \item T cells
    \end{choices}
    
    \item After meiosis-I daughter cells will have:
    \begin{choices}
      \item twice as much DNA as a haploid gamete
      \item as much as a haploid gamete
      \item four times as much as a haploid gamete
      \item None of the above
    \end{choices}
    
    \item Resolving power of a microscope is a function of:
    \begin{choices}
      \item wavelength of light used
      \item numerical aperture lens system
      \item refractive index
      \item wavelength of light used and numerical aperture of lens system
    \end{choices}
  \end{parts}
  
  \item State whether the following statements are True or False with brief justification:\pts{$3 \times 2 = 6$}
  \begin{parts}
    \item Majority of human cancers prevalent worldwide are caused due to viruses.
    \item Polytene chromosomes are found in amphibian oocytes and are giant meiotic prophase chromosomes.
    \item Glycosylation of lipids and proteins takes place in endosomes and nucleus, respectively.
  \end{parts}
  
  \item Define the following:\pts{$3 \times 1 = 3$}
  \begin{parts}
    \item RB gene
    \item Rhopheocytosis
    \item NOR.
  \end{parts}
\end{parts}

\pagebreak

\medskip\noindent\textbf{Question 2.} Answer the following:
\begin{parts}
  \item Describe the basic structure of immunoglobulin with suitable diagram. Discuss the characteristic features and functions of different types of immunoglobulins.\pts{6}
  \item Explain the role of Golgi complex in sorting of lysosomal enzymes.\pts{6}
\end{parts}

\medskip\noindent\textbf{Question 3.} Answer the following:
\begin{parts}
  \item Explain the various steps of the ribosome biogenesis in detail using well-labelled flow chart diagram.\pts{6}
  \item What is the role of the inner mitochondrial membrane in ATP synthesis?\pts{6}
\end{parts}

\medskip\noindent\textbf{Question 4.} Write detailed notes on the following:\pts{$3 \times 4 = 12$}
\begin{parts}
  \item Synaptonemal complex and its role in chromosome pairing
  \item Process of transport across plasma membrane
  \item Lampbrush chromosomes and their functions.
\end{parts}

\bigskip
\begin{center}
  \textbf{SECTION -- B} \\
  \textbf{(Elementary Biochemistry) (Marks: 35)}
\end{center}

\noindent\textbf{Question 1.} Answer the following parts:
\begin{parts}
  \item State whether the following statements are True or False and give brief justification in support of your answer:\pts{$3 \times 2 = 6$}
  \begin{parts}
    \item Cerebrosides are derivatives of sphingolipids.
    \item Shine-Dalgarno sequence helps in DNA-replication.
    \item DNA Pol activity is primer dependent.
  \end{parts}
  
  \item Define the following:\pts{$3 \times 1 = 3$}
  \begin{parts}
    \item Vector
    \item Epimer
    \item Helicase.
  \end{parts}
  
  \item Draw the structure of the following:\pts{$2 \times 1 = 2$}
  \begin{parts}
    \item A fatty acid with $18:3 \; (\Delta^{9, 12, 15})$ composition
    \item 2'-deoxycytidine.
  \end{parts}
\end{parts}

\medskip\noindent\textbf{Question 2.} Elucidate the following in diagrammatic form only:\pts{$2 \times 6 = 12$}
\begin{parts}
  \item Activation energy-based concept of enzyme catalysis
  \item Four-step events of chain elongation during protein synthesis in prokaryotes.
\end{parts}

\medskip\noindent\textbf{Question 3.} Answer the following:
\begin{parts}
  \item Discuss the characteristic properties of peptide bond.\pts{5}
  \item Describe the transcription initiation assembly in prokaryotes.\pts{7}
\end{parts}

\medskip\noindent\textbf{Question 4.} Answer the following:
\begin{parts}
  \item Discuss the different types of intra-molecular interaction which stabilizes 3D structure of a protein.\pts{6}
  \item What are disaccharides? Describe the structures and functions of any three important disaccharides.\pts{6}
\end{parts}

\vfill
\rule{\linewidth}{0.4pt}
\begin{center}\small
\href{https://bhu-exam.vercel.app/}{Click Here}\\[4pt]\href{https://me-aryan.vercel.app/}{Click Here}\\[4pt]\href{https://quashan-me.vercel.app/}{Click Here}
\end{center}

\end{document}
